\section{Conclusion}
\label{sec:conclusion}

We presented MultiEMiT, a multimodal dataset of Italian Social TV commentary annotated across five pragmatic layers: Plutchik's emotions, irony and sarcasm, offensiveness, target orientation, and semiotic relations. Exploring the dataset with compact Dual-Encoder models and Pixtral 12B on a shared test split highlighted three main takeaways:
\begin{itemize}
    \item \textbf{RQ1 (Visual Context):} The usefulness of images depends strongly on the task. For irony detection, visual features provide a clear and consistent improvement over text alone, because multimodal irony often comes from the contrast between the image and the text. For emotion recognition, images alone offer little benefit and help only when combined with auxiliary task labels. For semiotic relations, simply concatenating frozen text and image features does not improve over text alone.
    \item \textbf{RQ2 (Correlation Prompting):} The auxiliary annotation layers carry valuable information across tasks. Providing them in the input improves text-only emotion and irony detection, multimodal emotion recognition alongside images, and semiotic relation classification. However, on multimodal irony, adding auxiliary labels reduces the advantage gained from the visual features. For zero-shot Pixtral, adding correlation rules to the prompt provides a major boost in recognizing irony.
    \item \textbf{RQ3 (Architectural Paradigms):} Compact Dual-Encoders perform best on emotion and irony detection, maintaining more balanced predictions across rare classes. In contrast, fine-tuned Pixtral achieves the best overall result on semiotic relations, showing that classifying how text and image interact benefits from joint vision--language modeling. At the same time, Pixtral struggles with the minority Parallel class and tends to over-generate labels in multi-label emotion detection, requiring output truncation to the first two predicted emotions.
\end{itemize}

Looking ahead, our findings open several avenues for future work. First, expanding MultiEMiT beyond political talk shows (\textit{Propaganda Live}) and reality television (\textit{Pechino Express}) to scripted drama, live sports, and cultural broadcasts would help reduce genre-specific class skew and test cross-genre generalization, while restoring instance-level broadcast provenance would enable per-genre analyses. Second, the gap between frozen feature concatenation and joint cross-modal reasoning observed on semiotic relations motivates evaluating jointly pre-trained or end-to-end fine-tuned multimodal encoders \cite{radford2021clip}, as well as non-oracle correlation prompting in which auxiliary layers are predicted rather than given. Third, future studies could investigate affective annotation schemas tailored to Social TV dynamics, targeting engagement-driven states (\textit{hype}, \textit{amusement}, \textit{frustration}, \textit{empathy}, \textit{boredom}) alongside viewer stances (\textit{support}, \textit{opposition}, \textit{hatewatching}, \textit{neutral}) to better capture digital television audience reactions.